% Generated. Do not edit.
% Projection scriptural-date-location/evidence, version 1: the Scriptural Date and
% Location page, generated whole from research/chronology.toml, the Scripture
% chronology corpus and the source library by tools/tpt document-types refresh.
% document: liturgy/roman-rite/1962/propers/temporal/53-thirteenth-after-pentecost
\section*{Appendix: Evidence for the Scriptural Dates}
\label{sec:date-location-evidence}
{\small For each passage, as the Scriptural Date and Location table names it: what is dated, how the corpus weighs each answer, the source's own words, and the sources cited. A subject the corpus does not date is reported as such and given no date.\par}
{\small
\par\smallskip\noindent\textit{Offertory: Ps. 30:15--16 (Vg.; Heb. 31)}\par\nobreak
\begin{itemize}[leftmargin=1.2em,itemsep=1pt,topsep=1pt,parsep=0pt]
\item \textbf{Historical setting}. No date: the corpus holds a date for this that it reports from another edition by exception, which is not printed until the corpus records the source's own qualification of it (research pending).
\item \textbf{Historical setting} (preferred). David's flight from the face of Absalom: When he fled from the face of his son Absalom (The Holy Bible, Douay-Rheims (Challoner revision), Project Gutenberg eBook 1581: Ps. 3:1).
\item \textbf{Composition} (preferred). Psalms: Before c.~165 B.C. (``before the Maccabean period, around 165 B.C.; no individual psalm can be dated securely''; Catholic critical chronology; New American Bible, Revised Edition --- USCCB web state (2026), Psalms introduction).
\end{itemize}
\par\smallskip\noindent\textit{Introit: Ps. 73:20 (Vg.; Heb. 74), Ps. 73:19 (Vg.; Heb. 74), Ps. 73:23 (Vg.; Heb. 74), Ps. 73:1 (Vg.; Heb. 74)}\par\nobreak
\begin{itemize}[leftmargin=1.2em,itemsep=1pt,topsep=1pt,parsep=0pt]
\item \textbf{Historical setting}. No date: the chronology corpus has not yet dated this (research pending).
\item \textbf{Composition}. As under Offertory: Ps. 30:15--16 (Vg.; Heb. 31).
\end{itemize}
\par\smallskip\noindent\textit{Gradual: Ps. 73:20 (Vg.; Heb. 74), Ps. 73:19 (Vg.; Heb. 74), Ps. 73:22 (Vg.; Heb. 74)}\par\nobreak
\begin{itemize}[leftmargin=1.2em,itemsep=1pt,topsep=1pt,parsep=0pt]
\item \textbf{Historical setting}. As under Introit: Ps. 73:20 (Vg.; Heb. 74), Ps. 73:19 (Vg.; Heb. 74), Ps. 73:23 (Vg.; Heb. 74), Ps. 73:1 (Vg.; Heb. 74).
\item \textbf{Composition}. As under Offertory: Ps. 30:15--16 (Vg.; Heb. 31).
\end{itemize}
\par\smallskip\noindent\textit{Alleluia: Ps. 89:1 (Vg.; Heb. 90)}\par\nobreak
\begin{itemize}[leftmargin=1.2em,itemsep=1pt,topsep=1pt,parsep=0pt]
\item \textbf{Historical setting}. As under Introit: Ps. 73:20 (Vg.; Heb. 74), Ps. 73:19 (Vg.; Heb. 74), Ps. 73:23 (Vg.; Heb. 74), Ps. 73:1 (Vg.; Heb. 74).
\item \textbf{Composition}. As under Offertory: Ps. 30:15--16 (Vg.; Heb. 31).
\end{itemize}
\par\smallskip\noindent\textit{Communion: Wis. 16:20}\par\nobreak
\begin{itemize}[leftmargin=1.2em,itemsep=1pt,topsep=1pt,parsep=0pt]
\item \textbf{Composition} (disputed). The Book of Wisdom: In or after the time of Ptolemy IV Philopator (221--204 B.C.) (``the time of ... Ptolemy IV Philopator (221-204 B.C.)''; The Catholic Encyclopedia, Volume 15 (1912)); In or after the time of Ptolemy VII Physicon (145--117 B.C.) (``the time of ... Ptolemy VII Physicon (145-117 B.C.)''; The Catholic Encyclopedia, Volume 15 (1912)). Place: Alexandria in Gigot's account (The Catholic Encyclopedia, Volume 15 (1912)).
\end{itemize}
\par\smallskip\noindent\textit{Gospel: Luke 17:11--19}\par\nobreak
\begin{itemize}[leftmargin=1.2em,itemsep=1pt,topsep=1pt,parsep=0pt]
\item \textbf{Narrated event}. No date: the chronology corpus has not yet dated this (research pending).
\item \textbf{Composition} (disputed). The Gospel of St Luke: c.~A.D.~70 (``About the year 70''; The Catholic Encyclopedia, Volume 14 (New York, 1912)); Before the end of the Roman imprisonment, when the Acts was finished (De auctore, de tempore compositionis et de historica veritate Evangeliorum secundum Marcum et secundum Lucam (1912), responsa-i-ix).
\end{itemize}
\par\smallskip\noindent\textit{Epistle: Gal. 3:16--22}\par\nobreak
\begin{itemize}[leftmargin=1.2em,itemsep=1pt,topsep=1pt,parsep=0pt]
\item \textbf{Composition} (answers differ). The Epistle to the Galatians: preferred: A.D.~58 (``anno Christi 58''; Commentaria in omnes D. Pauli epistolas (1614), prooemium; Commentaria in omnes D. Pauli epistolas (1614), epistola-ad-galatas-argumentum); alternative: A.D.~49--50 (``A.D. 49-50''; The Catholic Encyclopedia, Volume 6 (1909)); alternative: c.~A.D.~53--54 (``about A.D. 53, 54''; The Catholic Encyclopedia, Volume 6 (1909)); alternative: A.D.~56 (``(1 and 2 Corinthians; Galatians), 56''; The Catholic Encyclopedia, Volume 11 (New York, 1911)); alternative: A.D.~57--58 (``A.D. 57 or 58''; The Catholic Encyclopedia, Volume 6 (1909)).
\end{itemize}
\par}
